\chapter{Lecture 12}

\begin{definition}[Group algebra]
    Let \(G\) be a finite group, and let \(R\) be a commutative unital ring.
    The group algebra \(RG\) is the free \(R\)-module with basis \(G\), whose
    multiplication is determined by the group multiplication of \(G\).
    
    Thus every element of \(RG\) can be written uniquely in the form
    \[
        x=\sum_{s\in G} x_s s,
        \qquad
        y=\sum_{t\in G} y_t t,
    \]
    where \(x_s,y_t\in R\). Their product is given by
    \[
        xy
        =
        \sum_{s,t\in G} x_s y_t st
        =
        \sum_{g\in G}
        \left(
            \sum_{\substack{s,t\in G\\ st=g}} x_s y_t
        \right) g.
    \]
    \end{definition}
    
    From now on, let \(k\) be a field, and consider the group algebra \(kG\).
    
    \begin{definition}
    Define a \(k\)-linear form
    \[
        \pi:kG\longrightarrow k
    \]
    by
    \[
        \pi(g)=
        \begin{cases}
            1, & g=1,\\
            0, & g\neq 1,
        \end{cases}
    \]
    for all \(g\in G\). Equivalently, \(\pi(x)\) is the coefficient of the identity
    element \(1\in G\) in \(x\in kG\).
    \end{definition}
    
    \begin{lemma}
    Let \(x,y\in kG\). Then
    \[
        \pi(xy)=\pi(yx).
    \]
    \end{lemma}
    
    \begin{proof}
    Write
    \[
        x=\sum_{s\in G}x_s s,
        \qquad
        y=\sum_{t\in G}y_t t.
    \]
    Then the coefficient of \(1\in G\) in \(xy\) is
    \[
        \pi(xy)
        =
        \sum_{\substack{s,t\in G\\ st=1}}x_s y_t
        =
        \sum_{s\in G}x_s y_{s^{-1}}.
    \]
    Similarly,
    \[
        \pi(yx)
        =
        \sum_{\substack{t,s\in G\\ ts=1}}y_t x_s
        =
        \sum_{t\in G}y_t x_{t^{-1}}.
    \]
    After replacing \(t\) by \(s^{-1}\), we get
    \[
        \pi(yx)
        =
        \sum_{s\in G}y_{s^{-1}}x_s
        =
        \sum_{s\in G}x_s y_{s^{-1}}
        =
        \pi(xy),
    \]
    because \(k\) is commutative.
    \end{proof}
    
    \begin{lemma}
    The map
    \[
        \psi:kG\longrightarrow \operatorname{Hom}_k(kG,k),
        \qquad
        x\longmapsto \bigl(y\mapsto \pi(xy)\bigr),
    \]
    is a \(k\)-linear isomorphism.
    \end{lemma}
    
    \begin{proof}
    It is clear that \(\psi\) is \(k\)-linear.
    
    Let \(f\in \operatorname{Hom}_k(kG,k)\). We claim that the inverse image of
    \(f\) under \(\psi\) is
    \[
        \psi^{-1}(f)
        =
        \sum_{g\in G} f(g^{-1})g.
    \]
    Indeed, set
    \[
        x=\sum_{g\in G} f(g^{-1})g.
    \]
    For any \(h\in G\), we have
    \[
        \psi(x)(h)
        =
        \pi(xh)
        =
        \pi\left(
            \sum_{g\in G} f(g^{-1})gh
        \right).
    \]
    The coefficient of \(1\) occurs precisely when \(gh=1\), that is,
    \(g=h^{-1}\). Therefore
    \[
        \psi(x)(h)
        =
        f\bigl((h^{-1})^{-1}\bigr)
        =
        f(h).
    \]
    Since \(G\) is a basis of \(kG\), this proves \(\psi(x)=f\). Hence \(\psi\) is
    surjective. The same formula also shows uniqueness of \(x\), so \(\psi\) is
    injective. Therefore \(\psi\) is a \(k\)-linear isomorphism.
    \end{proof}
    
    \begin{definition}[Class function]
    A function
    \[
        f:G\longrightarrow k
    \]
    is called a class function if
    \[
        f(gxg^{-1})=f(x)
    \]
    for all \(g,x\in G\).
    \end{definition}
    
    \begin{remark}
    If \(f:G\to k\) is a class function, then
    \[
        f(xy)=f(yx)
    \]
    for all \(x,y\in G\). Indeed,
    \[
        xy=x(yx)x^{-1},
    \]
    so \(xy\) and \(yx\) are conjugate.
    
    Extending \(f\) linearly, we obtain a \(k\)-linear map
    \[
        f:kG\longrightarrow k.
    \]
    Thus class functions on \(G\) may be regarded as special linear forms on
    \(kG\).
    \end{remark}
    
    \begin{remark}
    There is a natural identification
    \[
        \{\,\text{functions }G\to k\,\}
        \cong
        \operatorname{Hom}_k(kG,k)
        =
        (kG)^*.
    \]
    Indeed, since \(G\) is a \(k\)-basis of \(kG\), a linear form on \(kG\) is
    uniquely determined by its values on the basis elements \(g\in G\).
    \end{remark}

    \begin{definition}[Symmetric algebra]
        Let \(A\) be a finite-dimensional \(k\)-algebra. We say that \(A\) is a
        \emph{symmetric algebra} if there exists a \(k\)-bilinear form
        \[
            \langle -,-\rangle:A\times A\longrightarrow k
        \]
        such that:
        
        \begin{enumerate}
            \item \(\langle -,-\rangle\) is nondegenerate;
            \item \(\langle a,b\rangle=\langle b,a\rangle\) for all \(a,b\in A\);
            \item \(\langle ab,c\rangle=\langle a,bc\rangle\) for all \(a,b,c\in A\).
        \end{enumerate}
    \end{definition}

    \begin{definition}[Equivalent definition]
        A finite-dimensional \(k\)-algebra \(A\) is called symmetric if
        \[
            A\cong A^*
        \]
        as \(A\)-bimodules, where
        \[
            A^*=\operatorname{Hom}_k(A,k).
        \]
    \end{definition}
    
    \begin{corollary}
    The group algebra \(kG\) is a symmetric algebra.
    \end{corollary}
    
    \begin{proof}
    The bilinear form
    \[
        \langle x,y\rangle
        =
        \pi(xy)
    \]
    is symmetric by the first lemma:
    \[
        \langle x,y\rangle
        =
        \pi(xy)
        =
        \pi(yx)
        =
        \langle y,x\rangle.
    \]
    The second lemma says precisely that the induced map
    \[
        kG\longrightarrow (kG)^*,
        \qquad
        x\longmapsto \langle x,-\rangle,
    \]
    is a \(k\)-linear isomorphism. Therefore the bilinear form is nondegenerate,
    and hence \(kG\) is a symmetric algebra.
    \end{proof} 

